\documentclass[
10pt, % Main document font size
a4paper, % Paper type, use 'letterpaper' for US Letter paper
oneside, % One page layout (no page indentation)
%twoside, % Two page layout (page indentation for binding and different headers)
headinclude,footinclude, % Extra spacing for the header and footer
BCOR5mm, % Binding correction
]{scrartcl}
\input{stru/structure.tex}
\hyphenation{Fortran hy-phen-ation} 

\title{\normalfont{FTML practical session 2}} 
\author{\spacedlowsmallcaps{}
} 
\date{\today}
\begin{document}
\renewcommand{\sectionmark}[1]{\markright{\spacedlowsmallcaps{#1}}}
\lehead{\mbox{\llap{\small\thepage\kern1em\color{halfgray} \vline}\color{halfgray}\hspace{0.5em}\rightmark\hfil}} 
\pagestyle{scrheadings}
\maketitle 
\setcounter{tocdepth}{3} 

\begin{figure}[htpb]
    \centering
    \includegraphics[width=0.8\linewidth]{ols_test_errors_10_repetitions}
    \label{fig:risks}
\end{figure}

\tableofcontents

\pagebreak
\section*{\large\color{Blue}Introduction}

The goal of this session is to keep exploring the concepts linked to risks: Bayes risk,
empirical risk, overfitting. We keep exploring them in a linear regression
context, but dive a bit deeper technically and start studying the influence of
the dimensions (number of samples, number of features) on the statistical
properties of the results. Some aspects of this session are a bit technical, it
is normal if you need to take some time to understand all the objects used in
the session.

There are some template files in \textbf{tp\textunderscore 02\textunderscore ols/} that you can use in the repo, but you can also
start from scratch with your own files if you prefer.

Inside \textbf{2\textunderscore 5\textunderscore overfitting/}
\begin{itemize}
    \item \textbf{main\textunderscore ols.py}
    \item \textbf{utils\textunderscore plots.py}
    \item \textbf{utils\textunderscore algo.py}
    \item \textbf{constants.py}
\end{itemize}

\section{\large\color{Blue}Linear regression}
\label{sec:ols}

A \textbf{{linear model}}, such as the Ordinary least squares (OLS, presented below), can be interpreted as
predicting an output value (dependent variable) from combining the contributions from the $d$
\textbf{{features}} of the input data (independent variables), in a linear
way. This can be useful for classification as well as for regression.

\subsection{A simple example of linear function}%

For instance, if I want to predict the amount of money that I will spend when
buying some clothes, I can use a linear mapping (aka function). We assume that there are $d$ possible clothes in a shop. If $\theta\in \mathbb{R}^d$ is a vector that contains the price
of each type of clothe, and $x\in\mathbb{N}^d$ the number of each type of clothe that I buy,
then I have to spend $x^T\theta$. If there exists $d=4$ types of clothes with a
price $\theta_i$:
\begin{itemize}
    \item socks: $\theta_1=2$
    \item T-shirts: $\theta_2=25$
    \item pants: $\theta_3 = 50$
    \item hats = $\theta_4 = 20$
\end{itemize}

\begin{equation*}
\theta=\begin{pmatrix}
2\\
25\\
50\\
20
\end{pmatrix}
\end{equation*}

If I want to buy $10$ socks, $2$ T-shirts, $1$ pants and $1$ hat, then $x^T=(10,
2, 1, 1)$ and I spend
\begin{equation}
    \begin{aligned}
        \label{eq:}
        x^T\theta &= 10\times 2 + 2\times 25 + 1\times 50 + 1\times 20\\
        &= 140
    \end{aligned}
\end{equation}

\subsection{Linear models}

When doing a \textbf{linear regression} in a supervised learning context, we are using a linear function of the form $x\mapsto x^T\theta$ in order to approximate the outputs, where $x\in \mathbb{R}^d$ is an element of our input space, and $\theta\in \mathbb{R}^d$. The objective is then to find the best $\theta$ for this purpose.
\\

Of course, not all phenomena can be approximated well in a linear way. For many tasks, neural networks and many other models outperform linear regression.  However,
linear regression is a foundation for more advanced modelisation that we will
study in future classes (feature maps, kernel methods, neural networks, etc).
\\

It is also important to note that linear regression can often be used as a first benchmark when comparing different kinds of models, in a supervised learning context. A model that is more difficult to train or takes a longer time to train than a linear regression should have a better performance than linear regression, in order to be worth using it !
\\

As we will study the influence of the dimensions of the problem $n$ (number of
samples) and $d$ (number of features of each sample), we
will work with abstract datasets, instead of a given "physical dataset".
However, you can always refer to the previous example in order to interpret
linear models that we will use.

\subsection{Formalization}
\label{sec:formalization}

Let us abstract the notations a little bit. In the OLS setting,
\begin{itemize}
	\item $ \mathcal{X} = \mathbb{R}^d$ (input space)
	\item $ \mathcal{Y}  = \mathbb{R} $ (output space)
\end{itemize}

An input sample is stored in a column vector $x_a\in \mathbb{R}^d$, with $a\in \mathbb{N}$ like
\begin{equation}
    x_a=
\begin{pmatrix}
x_{a1}\\
...\\
x_{ai}\\
...\\
x_{ad}
\end{pmatrix}\in \mathbb{R}^d
\end{equation}

The estimator is a \textbf{{linear mapping}} parametrized by $ \theta\in
\mathbb{R}^d$. The prediction associated with $x\in \mathbb{R}^d$ is $f(x)=
\theta^Tx = x^T\theta$. If we use the squared-loss, the discrepancy between two real numbers $y$ and $y'$ writes $l(y,y')=(y-y')^2$. Finally, the input dataset is stored in the \textbf{{design matrix}} $X\in \mathbb{R}^{n\times d}$.

\begin{equation}
X=
\begin{pmatrix}
x_{1}^T\\
...\\
x_{i}^T\\
...\\
x_{n}^T
\end{pmatrix}
=
\begin{pmatrix}
    x_{11}, \:\dots\:,\: x_{1j}\:, \:\dots\: x_{1d}\\
...\\
    x_{i1}, \:\dots\:,\: x_{ij}\:, \:\dots\: x_{id}\\
...\\
    x_{n1}, \:\dots\:,\: x_{nj}\:, \:\dots\: x_{nd}\\
\end{pmatrix}
\end{equation}

and the output labels are stored in a vector

\begin{equation}
    y=
\begin{pmatrix}
y_{1}\\
...\\
y_{i}\\
...\\
y_{n}
\end{pmatrix}\in \mathbb{R}^n
\end{equation}


\subsection{Empirical risk and generalization error}

The \textbf{empirical risk} of an estimator $\theta$ (when we talk about an estimator, it is here equivalent to refer to $\theta$ directly or to the mapping defined by $\theta$) writes:


\begin{equation}
    \begin{aligned}
	R_n(\theta) &= \frac{1}{n} \sum^{n}_{i=1} (y_i-f(x_i))^2\\
		    &= \frac{1}{n} \sum^{n}_{i=1} (y_i-\theta^Tx_i)^2\\
    \end{aligned}
\end{equation}

{\color{Purple}{How can we write the empirical risk using only matrices and vectors, $X$, $\theta$ and $y$ defined in \ref{sec:formalization} and an euclidean norm ?}}
\\

By contrast, the \textbf{generalization error} (risque réel) of an estimator does not depend on the dataset. It is a fixed number equal to:

\begin{equation}
    R(\theta) = E[(y-\theta^Tx)^2]
\end{equation}

where the expected value is taken over the law of $x\in \mathbb{R}^d$ and $y \in\mathbb{R}$, the random variables representing the input and output respectively. This law is unknown, in general.

\pagebreak

\section{\large\color{Blue}Ordinary least squares: empirical risk and overfitting}

In this exercice, we want to study the influence of the dimensions $n$ and $d$ on the amount of \textbf{overfitting} of a linear regression. In order to do so, we will need to make statistical assumptions about the data that we use.

\subsection{OLS estimator}

The \textbf{OLS estimator}, noted $\hat{\theta}$, is the value of $\theta$ that minimizes $R_n(\theta)$. It is thus the solution to the problem of \textbf{empirical risk minimization}, a standard approach in supervised learning. We admit the following proposition (that we will prove later in the course):


\begin{proposition}{Closed form solution}

    We keep the notations from \ref{sec:formalization}. If $X$ is injective, there exists a unique minimiser of $R_n(\theta)$, called the \textbf{{OLS estimator}}, given by

    \begin{equation}
	\label{def:ols}
	\hat{ \theta}= (X^TX)^{-1}X^Ty
    \end{equation}
\end{proposition}

%{\color{Purple}{With the results of the exercises made today, you can try to prove this result.}}

\subsection{Excess risk the OLS estimator}

Given an estimator $\theta$, we are interested in its \textbf{excess risk}. It is the
difference between its generalization error and the bayes risk. We will evaluate the excess risk of the OLS estimator as a function of $n$ and $d$, for a given statistical setting, in order to observe the results of figure \ref{fig:test_errors_10}, \ref{fig:test_errors_1000}, and
\ref{fig:test_errors_10000}.

\subsection{Statistical setting}%
\label{ssub:Statistical setting}

In order to compute these quantities, it is necessary to make statistical
assumptions. We will use the \textbf{{linear model}}, with \textbf{{fixed
design}}, a classical framework to analyze OLS. In this setting, we assume that
$X$ is given and fixed, and that there exists a vector $ \theta^*\in \mathbb{R}^d$, such that $\forall i\in \{1, \dots, n\}$, 
\begin{equation}
    y_i = x_i^T\theta^*+\epsilon_i
\end{equation}

where for all $i\in \{1, \dots, n\}$, $\epsilon_i$ are independent, with
expectation $E[\epsilon_i]=0$ and variance $E[\epsilon_i^2]=\sigma^2$. The
$\epsilon_i$ represent a variability in the output, that is due to
\textbf{{noise}}, or to the presence of unobserved variables. Put together in a
vector $\epsilon\in \mathbb{R}^n$, this allows to write

\begin{equation}
    y = X\theta^*+\epsilon
\end{equation}

{\color{Purple}{Considering the previous practical session observations, what do
we expect to be the Bayes estimator ?}}
\\

We admit for now (but you can also prove it) that the corresponding Bayes risk is $\sigma^2$.
\\

\subsection{Verification of the Bayes estimator}%
\label{sub:Verification of the Bayes estimator}


{\color{Purple}{For a given $n$ and $d$ (with $d\leq n$), run a simulation that reproduces the fixed design, linear model
statistical setting and verify that the the estimator that you have chosen
achieves the Bayes risk. Verify that different estimators have worse (larger)
risks. You need to generate first a $X$ and a $\theta^*$.}}
\\

You will again need to approximate the real risk by the test error, a
number of times that is sufficient to observe the convergence of the law of
large numbers.
\\

You should observe something similar to figure \ref{fig:convergence}.

\begin{figure}[htpb]
    \centering
    \includegraphics[width=0.8\textwidth]{convergence}
    \caption{Convergence of the test error to the generalization error.}
    \label{fig:convergence}
\end{figure}

\subsection{Empirical risk minimization and impact of the dimensions on the OLS overfitting}

We have seen that in this statistical setting, the Bayes risk is $\sigma^2$. In
practical applications, we of course do not have access to $\theta^*$ and we
must find a relevant $\theta$, based on the data only. We must resort to
\textbf{empirical risk minimization}.
\\

We want to compare the risk of the OLS estimator (obtained by empirical risk minimization) to the Bayes risk, to know how
good the OLS estimator compares to the Bayes estimator. Again, remember that in
practical applications, we will not have access to the Bayes estimator or to the
Bayes risk.
\\

{
    \color{Purple}{Run a simulation that approximates the expected value of the generalization error of the OLS
    estimator. Remember that as the train data should be seen as random
    variables, the OLS estimator $\hat{\theta}$ himself should also be seen as
    a random variable, and so should be its generalization error. This is why
    we are approximating an \textbf{expected value} !

    To do so, we can generate several train and test data, and average the
    results of the tests errors in order to approximate this exepcted value. Depending on the number of repetitions of the experiment, on which we perform
the average, we should observe results simlar to that of figures
\ref{fig:test_errors_10}, \ref{fig:test_errors_1000}, and
\ref{fig:test_errors_10000}.
}
}
\\

\textbf{Remarks}
\begin{itemize}
	\item In order to generate the $X$ matrices in various dimensions, you can use uniformly distributed entries. This should ensure that $X$ is injective. 
	\item We will prove during the lectures the different results mentioned in this
exercice.
\end{itemize}

\begin{figure}[H]
    \centering
    \includegraphics[width=1\linewidth]{ols_test_errors_10_repetitions}
    \caption{Dependence of the test error (which approximates the generalization error (risque réel)) to the dimensions.}
    \label{fig:test_errors_10}
\end{figure}

\begin{figure}[H]
    \centering
    \includegraphics[width=1\linewidth]{ols_test_errors_1000_repetitions}
    \caption{Dependence of the test error (which approximates the generalization error (risque réel)) to the dimensions.}
    \label{fig:test_errors_1000}
\end{figure}

\begin{figure}[H]
    \centering
    \includegraphics[width=1\linewidth]{ols_test_errors_10000_repetitions}
    \caption{Dependence of the test error (which approximates the generalization error (risque réel)) to the dimensions.}
    \label{fig:test_errors_10000}
\end{figure}

We can make the following, intuitive observations:
\begin{itemize}
	\item  overfitting is smaller when $n$ is larger
	\item  overfitting is larger when $d$ is larger
\end{itemize}

It is possible to show that the amount of overfitting is given by $\sigma^2\frac{d}{n}$. Let me know if you are interested in proving this result.


\subsection{When $d>n$}%
\label{sub:When $d>n$}


If $d>n$, the $\sigma^2d/n$ law is not respected !
\\

{\color{Purple}{Why can we still output an OLS estimator when $d > n$ ?}}
\\

\begin{figure}[H]
    \centering
    \includegraphics[width=0.8\textwidth]{ols_risks_d_big}
    \caption{Test errors with cases where $d>n$.}
    \label{fig:ols_risks_d_big}
\end{figure}

{\color{Purple}{Optional: monitor the stability of the OLS estimator as a
function of $n$ and $d$, for instance by computing the sample standard deviation of
the random variable  $||\hat{\theta}-\theta^*||/||\theta^*||$ as a function of
$n$ and $d$. When $d$ is closer to $n$, $\hat{\theta}$ can become very
unstable (a small variation of the noise might lead to a large variation of
$\hat{\theta}$). This behavior depends on the conditioning of $X^TX$, as we will
see later.)}}

\begin{figure}[H]
    \centering
    \includegraphics[width=0.8\textwidth]{ols_stds}
    \caption{Instability of the OLS estimator when $d$ gets close to $n$.}
    \label{fig:}
\end{figure}

\end{document}
